\documentclass[10pt]{article}
\usepackage[T1]{fontenc}
\usepackage[utf8]{inputenc}
\usepackage[french]{babel}
\usepackage{amsmath,amssymb,mathrsfs,mathpazo}
\usepackage{geometry}
\geometry{a4paper,margin=2cm}
\usepackage{graphicx,array,tabularx,tabularray,colortbl,multirow,multicol,booktabs,enumitem,pifont,ulem,fancybox,diagbox,cancel}
\usepackage{pstricks,pstricks-add,pst-plot,pst-tree,pst-node,pst-text,pst-func,pst-math,pst-bspline,pst-3dplot,pst-eucl,pst-solides3d}
\let\pscalculate\undefined
\usepackage{pst-bezier}
\usepackage{tikz,pgfplots,tkz-tab}
\usetikzlibrary{arrows,arrows.meta,patterns,calc,positioning,trees,decorations.markings,shapes.geometric,shapes.misc,shapes.arrows,intersections,decorations.pathmorphing,decorations.pathreplacing,backgrounds,shadows,babel}
\pgfplotsset{compat=1.18}
\usepgfplotslibrary{fillbetween}
\usepackage[np]{numprint}
\usepackage{listings,makecell,pas-tableur,twoopt}
\let\barre\undefined
\usepackage[pstricks]{tablvar}
\definecolor{darkgreen}{RGB}{0,100,0}
\providecommand{\red}{\color{red}}
\providecommand{\blue}{\color{blue}}
\providecommand{\green}{\color{green}}
\providecommand{\black}{\color{black}}
\providecommand{\white}{\color{white}}
\providecommand{\gray}{\color{gray}}
\providecommand{\cyan}{\color{cyan}}
\providecommand{\magenta}{\color{magenta}}
\providecommand{\yellow}{\color{yellow}}
\providecommand{\violet}{\color{violet}}
\providecommand{\orange}{\color{orange}}
\providecommand{\euro}{\texteuro}
\providecommand{\EUR}{\texteuro}
\providecommand{\ieme}{\textsuperscript{e}}
\providecommand{\degres}{\ensuremath{^\circ}}
\providecommand{\degree}{\ensuremath{^\circ}}
\providecommand{\vv}[1]{\overrightarrow{#1}}
\providecommand{\Coord}[1]{#1}
\providecommand{\decosix}{}
\providecommand{\starredbullet}{\ensuremath{\bullet}}
\pagestyle{empty}
\setlength{\parindent}{0pt}
\usepackage{adjustbox}\DeclareUnicodeCharacter{2236}{:}\DeclareUnicodeCharacter{2190}{\ensuremath{\leftarrow}}
\providecommand{\R}{\mathbb{R}}
\providecommand{\N}{\mathbb{N}}
\providecommand{\Z}{\mathbb{Z}}
\providecommand{\Q}{\mathbb{Q}}
\providecommand{\D}{\mathbb{D}}
\providecommand{\C}{\mathbb{C}}
\providecommand{\e}{{\mathrm{e}}}
\providecommand{\vect}[1]{\overrightarrow{#1}}
\providecommand{\vectt}[1]{\overrightarrow{\mathrm{#1}}}
\providecommand{\vvt}[1]{\overrightarrow{\mathrm{#1}}}
\providecommand{\np}[2][]{{#2~#1}}
\providecommand{\vv}[1]{{\overrightarrow{#1}}}
\renewcommand{\d}{{\mathrm{d}}}
\providecommand{\degres}{{^\circ}}
\providecommand{\degre}{{^\circ}}
\providecommand{\blue}{\color{blue}}
\providecommand{\red}{\color{red}}
\providecommand{\enskip}{\quad}
\providecommand{\llbracket}{\lbrack\!\lbrack}
\providecommand{\rrbracket}{\rbrack\!\rbrack}
\providecommand{\euro}{{\text{€}}}
\providecommand{\no}{{\text{n°}}}
\providecommand{\interval}[2]{{[#1;#2]}}
\providecommand{\intervFO}[2]{{[#1;#2[}}
\providecommand{\dotfill}{\ldots}

\providecommand{\ds}{\displaystyle}
\providecommand{\textding}[1]{#1}
\providecommand{\Oij}{\(\mathrm{O};\vec{\imath},\vec{\jmath}\)}
\providecommand{\Oijk}{\(\mathrm{O};\vec{\imath},\vec{\jmath},\vec{k}\)}
\providecommand{\Ouv}{\(\mathrm{O};\vec{u},\vec{v}\)}
\providecommand{\intervFF}[2]{\ensuremath{[#1;#2]}}
\renewcommand{\euro}{\ifmmode\text{\texteuro}\else\texteuro\fi}
\DeclareUnicodeCharacter{00B0}{\ifmmode^{\circ}\else\textdegree\fi}
\DeclareUnicodeCharacter{20AC}{\ifmmode\text{\texteuro}\else\texteuro\fi}
% Énoncés : Éducation nationale. Transcription et corrigés : APMEP, Denis Vergès et contributeurs. Découpage et classement : CoopMaths, Sébastien Lozano, Éric Elter et contributeurs. Adaptation de la présentation : Maths Explorer.
% CC BY-SA : https://coopmaths.fr/www/about/
% https://coopmaths.fr/alea/static/e3c/2020/tex/e3c_2020_00_sujet33_4_cor.tex
\begin{document}
\noindent\textit{APMEP / CoopMaths. CC BY-SA. Présentation adaptée par Maths Explorer.}\par\medskip
\(g(x) = \e^x - x + 1\).

\subsection*{1.}

La fonction \(g\) est la somme de fonctions dérivables sur \(\mathbb{R}\), elle est donc dérivable sur \(\mathbb{R}\), et sa dérivée est :
\[
g'(x) = \e^x - 1.
\]

\subsection*{2.}

\begin{align*}
&\e^x - 1 > 0 \\
\iff &\e^x > 1 \\
\iff &x > 0.
\end{align*}
On en déduit que la fonction est :
\begin{itemize}[label=-]
    \item décroissante sur \([-5\,;\,0]\) de \(g(-5) = \e^{-5} - (-5) + 1 = \e^{-5} + 6\) à \(g(0) = 1 - 0 + 1 = 2\),
    \item croissante sur \([0\,;\,5]\) de \(g(0)\) à \(g(5) = \e^{5} - 5 + 1 = \e^{5} - 4\).
\end{itemize}

\subsection*{3.}

La question précédente a montré que le minimum de la fonction \(g\) sur l'intervalle \([-5\,;\,5]\) est \(2\), donc sur \([-5\,;\,5]\), \(g(x) \geqslant 2 > 0\).

\subsection*{4.}

On a, pour tout \(x \in [-5\,;\,5]\) :
\begin{align*}
f'(x) &= 1 + \dfrac{1 \times \e^x - x \times \e^x}{(\e^x)^2} \\
&= 1 + \dfrac{\e^x(1 - x)}{(\e^x)^2} \\
&= \dfrac{(\e^x)^2 + \e^x(1 - x)}{(\e^x)^2} \\
&= \dfrac{\e^x(e^x + 1 - x)}{(\e^x)^2} \\
&= \dfrac{e^x + 1 - x}{\e^x} \\
&= \dfrac{1}{\e^x} \times g(x).
\end{align*}

Le signe de \(f'(x)\) est donc celui de \(g(x)\). Or, on a vu que sur \([-5\,;\,5]\), \(g(x) > 0\), donc \(f'(x) > 0\) : la fonction \(f\) est donc strictement croissante de \(f(0) = 1\) à plus l'infini.

\subsection*{5.}

Si \(T_0\) est cette tangente, on sait qu'une équation de \(T_0\) est :
\[
M(x\,;\,y) \in T_0 \iff y - f(0) = f'(0)(x - 0).
\]
Avec \(f(0) = 1\) et \(f'(0) = \dfrac{g(0)}{\e^0}= g(0) = 2\), l'équation devient :
\begin{align*}
&M(x\,;\,y) \in T_0 \\
\iff &y - 1 = 2(x - 0) \\
\iff &y = 2x + 1.
\end{align*}

\end{document}
